\documentclass[11pt,a4paper]{article}

\usepackage[margin=0.88in]{geometry}
\usepackage[T1]{fontenc}
\usepackage[utf8]{inputenc}
\usepackage{lmodern}
\usepackage{amsmath,amssymb,amsthm,mathtools}
\usepackage{booktabs,tabularx,longtable,array}
\usepackage{enumitem}
\usepackage{xcolor}
\usepackage[numbers,sort&compress]{natbib}
\usepackage[colorlinks=true,linkcolor=blue!60!black,
 citecolor=blue!60!black,urlcolor=blue!65!black]{hyperref}
\usepackage[nameinlink,noabbrev]{cleveref}
\usepackage{microtype}

\setlength{\emergencystretch}{3em}
\def\Urlmuskip{0mu plus 1mu}
\newcolumntype{Y}{>{\raggedright\arraybackslash}X}

\hypersetup{
 pdftitle={TPC-MVP2: A Typed Global Route Decision after TPC-103--119},
 pdfauthor={Liang Wang},
 pdfsubject={Certified gates, stopped subroutes, open inputs, and the next minimal tests in the fixed-shift TPC program},
 pdfkeywords={fixed shift, typed dependency graph, proof audit, stop-go decision, endpoint budget, parity frontier}
}

\newtheorem{theorem}{Theorem}[section]
\newtheorem{proposition}[theorem]{Proposition}
\newtheorem{corollary}[theorem]{Corollary}
\newtheorem{lemma}[theorem]{Lemma}
\theoremstyle{definition}
\newtheorem{definition}[theorem]{Definition}
\theoremstyle{remark}
\newtheorem{remark}[theorem]{Remark}

\newcommand{\Lzero}{\mathrm{L0}}
\newcommand{\Lone}{\mathrm{L1}}
\newcommand{\Ltwo}{\mathrm{L2}}
\newcommand{\Lam}{\Lambda_{\mathrm{phys}}}
\newcommand{\epsend}{\frac1{400}}
\newcommand{\GO}{\textnormal{\textsc{go}}}
\newcommand{\RE}{\textnormal{\textsc{reroute}}}
\newcommand{\INF}{\textnormal{\textsc{infeasible}}}
\newcommand{\NT}{\textnormal{\textsc{not-testable}}}
\newcommand{\OPEN}{\textnormal{\textsc{open}}}
\newcommand{\PROVED}{\textnormal{\textsc{proved}}}

\title{\bfseries TPC-MVP2: A Typed Global Route Decision\\
after TPC-103--119\\[0.3em]
\large Certified Gates, Stopped Subroutes, Open Inputs,\\
and the Next Minimal Tests}
\author{Liang Wang\\
School of Artificial Intelligence and Automation\\
Huazhong University of Science and Technology, Wuhan 430074, P.R. China\\
\texttt{wangliang.f@gmail.com}}
\date{July 2026}

\begin{document}
\maketitle

\begin{abstract}
TPC-MVP1 resolved a proposed fixed-shift route into nine explicit
working hypotheses H1--H9.  TPC-103--119 then tested its positive
resonance, generic signed, determinant--zero, physical reconstruction,
endpoint, and reconnection interfaces.  We now place those results in
one typed directed acyclic graph.  Each node records its literal
object, fixed-shift quantifiers, normalization, evidence level, and
status.  This prevents a finite identity, a growing-family estimate,
and a fixed-\(h_0\) arithmetic saving from being interchanged.

The audit contains genuine progress.  It proves, among other
interfaces, an opened-atom cap, exact mass disintegration, a
provenance-preserving affine quotient, signed-filter identities,
post-bin determinant geometry, zero-mode duality, an exponent cone,
and exact tests for physical frames, localization, cover, endpoint
optimization, and native-atom reconnection.  It also stops several
specific subroutes: coefficient-blind distinctness, free-phase
prime-square arguments, and phase-blind complete-tail closure cannot
by themselves deliver the required conclusion.

The complete route nevertheless has neither a proved H1--H9
conjunction nor a theorem contradicting every allowed realization.
Actual growing inputs are missing at H2--H7, the executable physical
archives selected for the H6/H8 audits are absent, no alternative
complete H8 intertwining is present in the audited snapshot, and the
endpoint ledger is incomplete.  Therefore the exact global verdict is
\[
 \boxed{\text{\normalfont\textsc{not-testable}}}.
\]
It is not \(\GO\), not a global \(\INF\) verdict, and not evidence for
a prime-pair lower bound.  Local \(\RE\) labels remain attached only
to the subroutes actually stopped.  We finish with four minimal next
tests: two determinant--zero exponent certificates, a complete
native-atom archive, and one literal H3 block test.  This is an
\(\Lzero/\Lone\) route audit; no new \(\Ltwo\) cancellation theorem is
claimed.
\end{abstract}

\medskip
\noindent\textbf{Keywords:} fixed shift; typed proof graph; route
audit; endpoint budget; determinant zero mode; physical
reconnection.

\section{Scope and decision vocabulary}

Fix a prescribed nonzero shift \(h_0\), a smooth compactly supported
weight \(W\), and the original scalar packet
\(B_{h_0,\delta}(X)\) of TPC-15 \citep{WangTPC15}.  TPC-MVP1 is a
conditional architecture: its hypotheses H1--H9 were designed so
that their conjunction, in one literal normalization, implies
\[
 B_{h_0,\delta}(X)=o(X),
\]
which TPC-15 reconnects to the corresponding weighted fixed-shift
prime-pair asymptotic \citep{WangMVP1}.  The converse is not claimed;
the architecture may be stronger than its scalar target.

The present paper asks a narrower administrative question:
\begin{quote}
After TPC-103--119, can every required edge of that particular
architecture be called with theorem-backed data?
\end{quote}
It does not infer unproved arithmetic from the usefulness of the
architecture.

\begin{definition}[Evidence level]
\label{def:levels}
An \(\Lzero\) result is an abstract finite or functional-analytic
identity.  An \(\Lone\) result identifies that identity with the
literal TPC object and its normalization.  An \(\Ltwo\) result is a
growing, fixed-\(h_0\), coefficient-specific arithmetic estimate
with the power saving required by a hypothesis.  Passing from
\(\Lzero\) to \(\Lone\), or from \(\Lone\) to \(\Ltwo\), requires an
explicit theorem.
\end{definition}

\begin{definition}[Node record]
\label{def:record}
A node record is the six-tuple
\[
 {\mathfrak n}
 =
 (\text{object},\text{quantifiers},\text{normalization},
  \text{provenance},\text{evidence},\text{status}).
\]
It is \emph{typed} only when none of these fields is inferred from an
ambient, averaged-shift, surrogate, or differently normalized
object.
\end{definition}

\begin{definition}[Statuses and route verdicts]
\label{def:verdicts}
For an individual target, \(\PROVED\) means that its statement is a
theorem at the declared type; \(\OPEN\) means that the literal target
is specified but its estimate is unproved; \(\NT\) means that a
required archive, exponent, or compatibility input is absent, so the
current artifacts cannot execute the test.  A conditional interface
records an exact implication whose hypotheses are still open.

For route decisions, \(\GO\) requires every necessary node and the
strict endpoint.  \(\RE\) means that a named subroute is stopped while
an alternative remains.  \(\INF\) requires a theorem excluding every
allowed realization of the declared route.  A missing theorem is
neither \(\INF\) nor evidence against the underlying number-theoretic
statement.

The auxiliary label \(\textsc{stop route}\) means that all other gates
for one fully specified literal route are proved but its complete
endpoint registry has certified value at least \(1/400\).  It stops
that declared route only.  It becomes \(\RE\) when a typed alternative
is available, and it is not global \(\INF\) unless every allowed route
is excluded.
\end{definition}

\section{A typed decision calculus}

Let \({\cal H}=\{H1,\ldots,H9\}\) be the MVP1 gates.  An evidence
paper may supply an exact identity to a gate, isolate a missing
growing input, or stop a particular sufficient subroute.  None of
those actions changes the statement of another gate.

\begin{theorem}[Global decision rule]
\label{thm:decision}
For the declared MVP1 architecture, use the following mutually
exclusive rule.
\begin{enumerate}[label=(\roman*),leftmargin=2.2em]
\item Return \(\GO\) only if every gate in \({\cal H}\) has a proved
  typed record in one compatible quantifier scope and the actual
  physical ledger satisfies \(\Lam<1/400\).
\item Return \(\INF\) only if theorem-backed lower bounds or logical
  incompatibilities contradict every admissible realization of a
  necessary gate.
\item Return \(\RE\) for a selected subroute when that subroute is
  stopped but a typed alternative remains.
\item Otherwise, if a necessary estimate or typed input is open or
  absent, return \(\NT\).
\end{enumerate}
In particular, a stopped subroute cannot be promoted to global
\(\INF\), and an unknown endpoint token cannot be set to zero to
produce \(\GO\).
\end{theorem}

\begin{proof}
Item (i) is the invocation rule for the conditional synthesis theorem
of MVP1.  If one premise lacks a typed theorem, modus ponens is not
available.  Item (ii) requires universal coverage of the permitted
realizations; a counterexample to one sufficient bound or one
representation has weaker quantifiers.  Item (iii) preserves that
distinction.  The remaining case is precisely the absence of enough
information to execute either the upper or lower certificate, which
is \(\NT\) by \cref{def:verdicts}.  Finally, inserting zero for an
unknown nonnegative loss changes the physical optimization problem
rather than solving it.
\end{proof}

\subsection{The dependency graph}

The audit graph separates evidence bundles from gates and gates from
the final conditional synthesis:
\[
\begin{array}{c@{\quad}c@{\quad}c}
\boxed{\text{TPC-103--109}}&
\boxed{\text{TPC-110--112}}&
\boxed{\text{TPC-113--119}}\\
\downarrow&&\downarrow\qquad\downarrow\\[-0.2em]
\{H2,H3\}&\{H5\}&\{H1,H4,H6,H7,H8,H9\}\\
\multicolumn{3}{c}{\searrow\qquad\downarrow\qquad\swarrow}\\[-0.2em]
\multicolumn{3}{c}{
\boxed{\text{typed H1--H9 conjunction at fixed }h_0}}\\
\multicolumn{3}{c}{\downarrow}\\[-0.2em]
\multicolumn{3}{c}{
\boxed{B_{h_0,\delta}(X)=o(X)\ \text{within the MVP1 route}}.}
\end{array}
\]
This display is an audit DAG, not a claim that one evidence paper
proves the gate below it.  Every H-node is an ancestor of the final
conditional target; its status is given in \cref{tab:gates}.

\section{Paper-by-paper audit: TPC-103--119}

\Cref{tab:papers-a,tab:papers-b} record the strongest exact result and
the missing item that prevents overpromotion.  ``L1 interface'' can
include an exact finite identity on the literal carrier; it does not
mean that the required growing estimate follows.

\begin{table}[htbp]
\centering
\caption{Positive, generic, and determinant--zero audits.}
\label{tab:papers-a}
\scriptsize
\begin{tabularx}{\textwidth}{@{}p{0.055\textwidth}p{0.42\textwidth}
 p{0.16\textwidth}Y@{}}
\toprule
TPC & Exact contribution & Status & Unresolved or stopped item\\
\midrule
103 & Opened-atom cap \(W_X=X^{o(1)}\) on the normalized literal
carrier. & \(\PROVED\), L1 & The other H2 ledgers are not implied.\\
104 & Exact principal-mass disintegration before width grouping. &
L1 interface & Actual endpoint and inverse-length mass bounds are
\(\OPEN\).\\
105 & Provenance-preserving quotient by equal literal affine maps. &
\(\PROVED\), L1 & Growing class census \(U_X\) and distinct-map target
remain \(\OPEN\).\\
106 & Exact distinct-map character/Gram identity. &
L1; \(\RE\) & Geometry-only distinctness is stopped; actual profile
coherence remains \(\OPEN\).\\
107 & Exact signed-filter-before-majorant identity. &
\(\PROVED\), L1 & A growing signed-energy bound is \(\OPEN\).\\
108 & Exact literal-block \(TT^*\), aliasing, and obstruction
identities. & L1 interface & The H3 power saving is \(\OPEN\); a broad
Chowla-strength replacement is not available.\\
109 & Exact coherent prime-square free-phase interval and sharp
saturation. & L1; \(\RE\) & Coefficient-blind/free-phase route is
stopped; literal coefficient-specific phases remain \(\OPEN\).\\
110 & Post-bin SVD and
\(D_X=E_X\alpha_X\overline m_X\). &
\(\PROVED\), L1 identity & An actual lower bound for its three
growing inputs is \(\OPEN\).\\
111 & Exact zero-mode coarsening invariance and signed-prefix/BV
duality. & \(\PROVED\), L1 identity & The growing zero-mode exponent
\(\eta_Z\) is \(\OPEN\).\\
112 & Exact compatibility cone
\(\lambda_D\le2\eta_Z\) and separate-reserve rule. &
Conditional L1 & Neither actual exponent is certified; H5 is not
thereby proved.\\
\bottomrule
\end{tabularx}
\end{table}

\begin{table}[htbp]
\centering
\caption{Physical, endpoint, and reconnection audits.}
\label{tab:papers-b}
\scriptsize
\begin{tabularx}{\textwidth}{@{}p{0.055\textwidth}p{0.42\textwidth}
 p{0.16\textwidth}Y@{}}
\toprule
TPC & Exact contribution & Status & Unresolved or stopped item\\
\midrule
113 & Declared quotient condition number
\(\kappa_{\rm q}=\sigma_{\max}/\sigma_{\min}^{+}\) and a sharp
coherence obstruction. & Conditional L1 interface & The actual H9
operator factorization, norm conversion, and growing exponent are
\(\OPEN\).\\
114 & Literal residual coarea, native maps, Gram dictionary, and
collision-kernel obstruction. & L1 interface & Actual grouping
census, protected complement, H9 norm conversion, and native lower
transfer are \(\OPEN\).\\
115 & Exact fixed-\(h_0\) evaluation/Christoffel cost
\(K=b_0^*G^\dagger b_0\). & L1 interface & Actual profile subspace
and growing localization cost are \(\OPEN\).\\
116 & Exact conditional closure theorem and phase-blind obstruction. &
L1; local \(\RE\) & Complete outer tail bound is \(\OPEN\);
phase-blind closure is stopped.\\
117 & Exact full-cover test
\(P=BB^\dagger\), residual \((I-P)w\). &
\(\NT\) for H6 & The actual growing \(B,w\) archive and residual
estimate are absent.\\
118 & Nonduplicated physical ledger and exact rational primal--dual
endpoint logic. & \textsc{incomplete} & Several actual growing costs
are unknown; no strict endpoint certificate exists.\\
119 & Native-atom conservation and a sufficient canonical
reconnection-matrix criterion. & \(\NT\) for its H8 audit & The dated
snapshot has neither the selected complete leaf archive nor another
complete exact intertwining joining every TPC-15 atom to the final
synthesis.\\
\bottomrule
\end{tabularx}
\end{table}

The detailed statements appear in TPC-103--112
\citep{WangTPC103,WangTPC104,WangTPC105,WangTPC106,WangTPC107,
WangTPC108,WangTPC109,WangTPC110,WangTPC111,WangTPC112} and
TPC-113--119
\citep{WangTPC113,WangTPC114,WangTPC115,WangTPC116,WangTPC117,
WangTPC118,WangTPC119}.  TPC-18 itself treats a complete global block
cover as an additional hypothesis \citep{WangTPC18}; local block
identities therefore cannot silently supply the executable
certificates missing in TPC-117 or TPC-119.

\section{Projection onto H1--H9}

\begin{table}[htbp]
\centering
\caption{Current status of every required MVP1 gate.}
\label{tab:gates}
\small
\begin{tabularx}{\textwidth}{@{}p{0.07\textwidth}p{0.27\textwidth}
 p{0.16\textwidth}Y@{}}
\toprule
Gate & Literal role & Status & Reason\\
\midrule
H1 & Canonical physical synthesis & \(\NT\) & Exact frame and native
Gram criteria exist, but no complete actual-carrier frame/archive or
proved growing cost is supplied.\\
H2 & Literal positive resonance & \(\OPEN\) & The atom cap is proved;
principal/cross-map/coherence ledgers are not.  Some geometry-only
subroutes are stopped.\\
H3 & Restricted signed affine arithmetic & \(\OPEN\) & The literal
\(TT^*\) object is isolated, but no fixed-\(h_0\)
\(X^{-1/200+o(1)}\) correlation saving is proved.\\
H4 & Complete high/ultra return & \(\OPEN\) & Conditional closure and
an obstruction are known; the complete original-normalization outer
bound is not.\\
H5 & Determinant/zero compatibility & \(\NT\) & The cone is exact,
but the actual \(\lambda_D,\eta_Z\) inputs are absent.\\
H6 & Full-block physical return & \(\NT\) & The exact projector test
cannot be run without the growing physical \(B,w\) archive.\\
H7 & Prescribed-\(h_0\) localization & \(\OPEN\) & The exact
Christoffel cost is known; the literal profile and subcritical
growing bound are not.\\
H8 & Original-packet reconnection & \(\NT\) & Native conservation is
exact, but the sufficient proof-carrying leaf matrix selected in
TPC-119 is absent and the audited snapshot contains no alternative
complete exact intertwining.\\
H9 & Strict endpoint & \(\NT\) & The LP logic is exact; the actual
token registry is incomplete, so \(\Lam<1/400\) is uncertified.\\
\bottomrule
\end{tabularx}
\end{table}

\begin{theorem}[MVP2 global verdict]
\label{thm:verdict}
For the artifacts TPC-103--119 and the H1--H9 architecture stated in
TPC-MVP1,
\[
 \boxed{\operatorname{Verdict}_{103:119}=\NT.}
\]
There is no established \(\Ltwo\) gate in this audit.  The result is
not a global \(\INF\) certificate.
\end{theorem}

\begin{proof}
\Cref{tab:gates} contains four explicitly open gates and five gates
whose actual test lacks a required typed input.  Hence the premise of
\cref{thm:decision}(i) fails.  TPC-106, TPC-109, and TPC-116 stop
only coefficient-blind, free-phase, or phase-blind sufficient
subroutes; coefficient-specific signed routes remain logically
available.  Thus the universal premise of
\cref{thm:decision}(ii) is absent.  Applying item (iv) gives
\(\NT\).  None of the papers listed in
\cref{tab:papers-a,tab:papers-b} supplies the fixed-shift growing
cancellation demanded in \cref{def:levels}, so no \(\Ltwo\) gate is
established.
\end{proof}

\begin{remark}[What has moved forward]
The verdict does not mean ``no progress.''  Before these audits, an
apparent failure could be hidden among normalization, aliasing,
coherence, zero mode, cover, or reconnection.  The current graph gives
exact tests for those interfaces and identifies which data must be
measured or proved.  That is structural \(\Lone\) progress.  It is
not fixed-\(h_0\) parity cancellation.
\end{remark}

\section{The strict endpoint is an independent stop gate}

TPC-118 represents every physical reconstruction event by one
amplitude-loss token, merges exact duplicates, and replaces
constituent tokens only when a joint theorem actually subsumes them
\citep{WangTPC118}.  On a fixed affine route cell the decision is an
exact rational linear program with optimum \(\Lambda_*\).  Its
interpretation is
\[
\begin{array}{ccl}
\Lambda_*<1/400
 &\Longrightarrow& \GO\ \text{for a complete theorem-backed ledger},\\
\Lambda_*=1/400
 &\Longrightarrow& \text{equality stop},\\
\Lambda_*>1/400
 &\Longrightarrow& \text{stop that route},\\
\text{required token unknown}
 &\Longrightarrow& \NT.
\end{array}
\]

\begin{proposition}[No zero-filling and no reserve reuse]
\label{prop:endpoint}
The current endpoint cannot be certified by either
\begin{enumerate}[label=(\alph*),leftmargin=2em]
\item assigning exponent \(0\) to an unknown frame, grouping,
localization, tail, or cover cost; or
\item spending slack from the determinant--zero condition
\(\lambda_D\le2\eta_Z\) as physical endpoint slack.
\end{enumerate}
Moreover, equality \(\Lam=1/400\) does not pass H9.
\end{proposition}

\begin{proof}
Operation (a) changes an incomplete physical ledger into a different
optimization problem.  Operation (b) merges two ledgers that
TPC-112 and MVP1 declare independent absent a new intertwining
theorem \citep{WangTPC112,WangMVP1}.  Finally, the synthesis starts
from a packet saving of \(1/400\); equality consumes the whole
fixed-power reserve and leaves only an \(X^{1+o(1)}\) upper bound,
which does not certify \(o(X)\).
\end{proof}

\section{The next four minimal tests}

The decision graph recommends small falsifiable gates rather than a
new global claim.

\paragraph{TPC-121: actual post-bin fiber phase/participation.}
TPC-110 gives
\[
 D_X=E_X\alpha_X\overline m_X.
\]
The next determinant test should preserve the actual post-bin fibers
and certify the diagonal energy scale, the coefficient-specific
phase/kernel angle, and effective fiber participation.  Success
would supply a candidate \(\lambda_D\); failure could stop that
coherence mechanism without stopping every H5 route.

\paragraph{TPC-122: signed-prefix/BV/content transfer.}
Use the exact Abel/BV duality of TPC-111 on the literal ordered
fibers.  A theorem must retain the arithmetic sign and fixed \(h_0\),
bound signed prefixes, charge the reconstruction BV envelope, and
control the content remainder.  Its output is a candidate
\(\eta_Z\), to be compared with \(\lambda_D/2\), not a generic
nonzero-frequency spectral gap.

\paragraph{TPC-123: complete native-atom archive.}
Independently of H5, instantiate the TPC-119 canonical matrix:
every opened atom \((\ell,k,d)\) must have one recoverable provenance
key, every retained or soft leaf must appear, and column conservation
must hold at the same \(h_0\) and normalization.  A failed audit
locates the first missing, duplicated, rescaled, or surrogate branch.

\paragraph{TPC-124: one literal H3 block.}
Select one actual long generic affine block and preserve its M\"obius
weights, \(r\)-dependent phase, affine character, reconstruction
multiplier, all prefixes, and prescribed shift.  Compute the exact
normalization and test a rigorous upper/lower interval for its signed
\(TT^*\) correlation.  A finite experiment can choose the next
analytic lemma, but only a growing uniform theorem would count as
\(\Ltwo\).

\begin{corollary}[Interpretation of the next queue]
TPC-121 and TPC-122 attack the two missing exponent inputs of H5.
TPC-123 attacks the H1/H6/H8 physical archive, while TPC-124 attacks
H3 directly.  Progress on one branch may be published independently,
but none substitutes for the others in the H1--H9 conjunction.
\end{corollary}

\section{Reproducible audit and claim boundary}

The standard-library script
\texttt{experiments/tpc120\_global\_route\_audit.py} encodes the
paper-to-gate edges and verifies:
\begin{enumerate}[leftmargin=2em]
\item all nine H-nodes occur and have evidence ancestors;
\item no TPC-103--119 record is labeled as a proved \(\Ltwo\) result;
\item open or absent required inputs return \(\NT\);
\item a local stopped subroute is not global \(\INF\);
\item strict-below, equality, strict-above, and incomplete endpoint
  examples have four different outcomes;
\item unknown costs are not converted into a zero-cost certificate;
  and
\item the next queue is exactly TPC-121--124.
\end{enumerate}
The generated JSON is a regression artifact for the decision logic.
It is not a computation of any prime or M\"obius correlation.

\paragraph{Established.}
The status calculus, typed dependency graph, exact endpoint decision
rule, and paper-to-gate audit are \(\Lzero/\Lone\).  The cited papers
contain the individual exact identities summarized above.

\paragraph{Not established.}
This paper proves no H1--H9 conjunction, no strict actual-carrier
\(\Lam<1/400\) bound, no \(B_{h_0,\delta}(X)=o(X)\) estimate, and no
new \(\Ltwo\) fixed-shift cancellation.  It does not break the parity
barrier, prove a Hardy--Littlewood asymptotic, produce a prime-pair
lower bound, or prove the twin-prime conjecture.  Its conclusion is a
route decision: the path still has typed alternatives, but the
current archive is insufficient to decide their success.

\begingroup
\small
\setlength{\bibsep}{2pt plus 0.3ex}
\sloppy
\bibliographystyle{plainnat}
\bibliography{references}
\endgroup

\end{document}
